\subsection{Generalization and a targeted intervention}
\label{sec:generalization}

\paragraph{The separation curve depends on the deployment and task.}
The two added deployments were tested on the same six-layout follow-up fixtures as the original deployments. Across inline bar, two spaces, 40-pixel gap, separate lines, and opposite corners, respectively, \texttt{gpt-5.4} scored 0/20, 0/20, 18/20, 10/20, 13/20. \texttt{gpt-5-mini} scored 0/20, 0/20, 10/20, 13/20, 13/20. Figure~\ref{fig:separation} retains the task split; the full three-criterion cell estimates are in the artifact. The original-layout 15/20 result for \texttt{gpt-5.4} is therefore evidence about that fixture, not a deployment-wide separation threshold.

\paragraph{Measured distance in familiar interface shells.}
Figure~\ref{fig:generalization}A reports delimitation error against measured pixel distance. Across the 360 trials per deployment, exact task success was \texttt{gpt-5.6-terra} 108/360, \texttt{Kimi-K2.6} 116/360, \texttt{gpt-5.4} 115/360, \texttt{gpt-5-mini} 41/360. Delimitation is scored from the final read answer or last attempted typed content; correct content with a failed submission is kept separate. The full estimates retain layout, task, deployment, pixel distance, and DOM separation. Because the DOM variants have the same visual appearance, their contrast is a negative control rather than a test of structural access by vision. These controlled shells broaden layout coverage but do not establish real-site generalization.

Increasing the gap from 8 to 96 pixels reduced delimitation errors by 89.2, 83.3, 79.2, and 54.2 percentage points for \texttt{gpt-5.6-terra}, \texttt{Kimi-K2.6}, \texttt{gpt-5.4}, and \texttt{gpt-5-mini}, respectively; each paired 95\% bootstrap interval excluded zero. Hidden-DOM contrasts ranged from $-2.2$ to $+3.3$ points, with all intervals including zero. Full intervals are in the artifact. Lower delimitation error alone does not imply higher task success: missing answers and failed submissions remain separate outcomes.

The natural-boundary controls use separate table columns, neighboring cards, and labeled form fields ($N=240$). Across 60 trials per deployment, exact success was \texttt{gpt-5.6-terra} 59/60, \texttt{Kimi-K2.6} 45/60, \texttt{gpt-5.4} 47/60, \texttt{gpt-5-mini} 33/60. No delimitation errors were observed in these 240 trials. Figure~\ref{fig:generalization}A places these controls at their measured mean pixel distance. They change component boundaries and labels together and therefore support a compound layout comparison, not an isolated causal effect of DOM distance.

\paragraph{Explicit target formatting.}
Figure~\ref{fig:generalization}B compares contemporaneous paired baseline and format-prompt trials. \texttt{gpt-5.6-terra} changed from 0/20 to 20/20 (+100 percentage points; paired bootstrap 95\% CI +100 to +100); \texttt{Kimi-K2.6} changed from 0/20 to 15/20 (+75 percentage points; paired bootstrap 95\% CI +60 to +90); \texttt{gpt-5.4} changed from 1/20 to 11/20 (+50 percentage points; paired bootstrap 95\% CI +35 to +65); \texttt{gpt-5-mini} changed from 0/20 to 11/20 (+55 percentage points; paired bootstrap 95\% CI +50 to +65). This tests an answer-format instruction, not an oracle region or DOM cross-check. The comparison uses its own matched baseline because the calibrated inline fixture differs in viewport and canvas width from the original 0/40 result.

\begin{figure}[t]
  \centering
  \includegraphics[width=\linewidth]{generalization-intervention.pdf}
  \caption{\textbf{Generalization by measured distance and a targeted prompt intervention.} A: delimitation-error rates in tables, cards, and forms, pooling tasks and the two visually matched DOM boundaries ($N=40$ per deployment--layout--gap point). Diamonds show separate labeled components at their measured mean distance ($N=20$); all four deployments overlap at zero observed error. B: exact task success for inline-canvas baseline (open) and explicit-format prompts (filled), with $N=20$ per deployment and arm. Cell-level Wilson intervals and paired-effect intervals are provided in the artifact.}
  \label{fig:generalization}
\end{figure}

\paragraph{Iframe depth after calibration and position matching.}
The calibrated rerun obtained all 560 allocated outcomes. Pooled over the seven conditions, read/copy success was \texttt{gpt-5.6-terra} 70/70 / 70/70, \texttt{Kimi-K2.6} 70/70 / 56/70, \texttt{gpt-5.4} 70/70 / 44/70, \texttt{gpt-5-mini} 70/70 / 10/70. Reading was 10/10 at every depth/origin cell for every deployment. Copying was 10/10 throughout for \texttt{gpt-5.6-terra}; the other deployments varied non-monotonically across depths. There is no consistent depth gradient in these outcomes. Appendix~\ref{app:calibrated} retains all task--deployment--depth cells; paired comparisons against depth zero are in the artifact. Since input screenshots are position-matched across depths, residual differences cannot be attributed to the old fixed screenshot-scale or absolute-position artifact. Ten paired repeats per cell limit sensitivity; a small or nonsignificant contrast is not evidence of equivalence. The original variance decomposition remains only a descriptive audit in Appendix~\ref{app:variance-audit}.
